%% kapa: exam=1 round=36 year=2025 period=1 form=A subject=economics subject_ko=경제학원론 seq=38 number=78
%% source: Q-Net 감정평가사 기출문제 (cst00309) artlSeq=5231525 file=raw/1008/5231525/2025년 제36회 감정평가사 1차시험 1교시 문제지 원본.pdf
%% conversion: kordoc 4.24.51 -> tools/gen1.py (auto); review=proofread-sonnet; flags=none
\begin{question}{78}
주조차익(seigniorage)과 인플레이션 조세(inflation tax)에 관한 설명으로 옳은 것은?
\begin{choices}
  \item 실물변수가 변화하지 않는 장기 균형 상황에서는 주조차익과 인플레이션 조세는 같다.
  \item 인플레이션이 발생하지 않으면, 주조차익은 0이다.
  \item 통화량이 감소하면, 주조차익은 증가한다.
  \item 물가상승률이 통화증가율보다 크면, 인플레이션 조세가 주조차익보다 작다.
  \item 통화증가율이 0이고 인플레이션이 커지면, 주조차익도 커진다.
\end{choices}
\end{question}
\answer{1}
